\documentclass[10pt,a4paper]{amsart}
\usepackage[margin=19mm]{geometry}
\usepackage{amsmath,amssymb,mathtools}
\usepackage[scr=rsfs,cal=euler]{mathalfa}
\usepackage{xfrac}
\usepackage[unicode,hidelinks,bookmarks=false]{hyperref}
\newcommand{\ie}{i.e.\ }
\newcommand{\cf}{cf.\ }
\newcommand{\resp}{respectively\ }
\newcommand{\Subsection}{\subsection}
\providecommand{\nobreakdash}{\nobreak\hbox{-}}
\setlength{\emergencystretch}{2em}
\multlinegap0pt
\usepackage{amssymb}
\usepackage{enumerate}
\usepackage[all]{xy}
\usepackage{xcolor}
\newtheorem{cor}{Corollary}
\newtheorem{lemma}{Lemma}
\newcommand{\C}{\mathcal{C}}
\newcommand{\Csurv}{\C^{\text{surv}}}
\newcommand{\Z}{\mathbb{Z}}
\newcommand{\diam}{\operatorname{diam}}
\newcommand{\dist}{d}
\title{Approximating stable translation lengths\\ on fine curve graphs}
\author{Federica Fanoni, Sebastian Hensel, Fr\'ed\'eric Le Roux}
\date{Source: arXiv:2601.03412v1 (CC BY 4.0)}
\begin{document}
\maketitle

\noindent\emph{Source and licence.} Approximating stable translation lengths\\ on fine curve graphs. Version arXiv:2601.03412v1 (CC BY 4.0), distributed by arXiv under the Creative Commons Attribution 4.0 International licence, http://creativecommons.org/licenses/by/4.0/, as recorded in the arXiv licence metadata for this exact version and shown on the abstract page https://arxiv.org/abs/2601.03412v1. Full licence notice: \texttt{licenses/arxiv-2601.03412-CC-BY-4.0.txt}.\par\smallskip

\noindent\textit{Editorial scope.} Counted unit: the article's lemma lem:geodesics-between-iterates (source0.tex lines 407--416). The article's complete original proof of it is delivered with it (source0.tex lines 418--473, one \texttt{proof} environment of 56 lines). No mathematics is written, repaired or re-derived: the statement and the proof are the article's own lines, in the article's own notation, and no retained line is substituted or re-flowed.  Retained uncounted as the source-local context the proof needs: lines 236--243.  Nothing was omitted from any retained span, so the package carries no omission record.  No retained line contains a citation, so the wrapper carries no bibliography and no bibliography entry.  No retained line contains a figure environment, an image-inclusion command or drawing code, so no image is delivered and the package declares no asset dependency.  Editorial formatting only, in addition: an AMS \texttt{amsart} preamble that restates the article's own declarations and macros used by the retained lines, namely the environments \texttt{cor}, \texttt{lemma} on separate counters, together with the standard packages those lines need.  The retained environments print in this document as Corollary 1, Lemma 1 in order of appearance, whereas the article prints the same items with the numbering of its own full document.  The complete native source of the article as distributed in the arXiv e-print for this exact version is bundled unchanged as \texttt{sources/192224/source0.tex}.
\begin{cor}[Bounded Quasigeodesic Image Theorem]\label{cor:qBGIT}
  Let \(S\) be a surface of complexity at least three. For every $K>0$
  there is a constant \(B\) depending only on \(S\) and $K$ such that
  the following holds: for every subsurface \(Z\) which is either an
  annulus or has complexity at least four, if \(c_0,\dots,c_n\) is a
  $K$--quasigeodesic path such that \(\pi_Z(c_i)\neq\emptyset\) for
  every \(i\), then \(\diam \bigcup_{i=0}^n\pi_Z(c_i)\leq B\).
\end{cor}

\begin{lemma}\label{lem:geodesics-between-iterates}
Let \(\varphi\) be a pseudo-Anosov mapping class and \(a_0\in\Csurv(S)\). 
There are constants \(A\), \(K\) and $B$, depending on \(\varphi\) and \(a_0\), such that:
\begin{enumerate}
\item  for every \(n\in\Z\) and for every subsurface \(Z\), if $\pi_Z(\varphi^n(a_0)) \neq \emptyset$ then
\[\dist_Z(\pi_Z(\varphi^n(a_0)),\pi_Z(\lambda^+))\leq A;\]
\item  any geodesic between \(\varphi^{-n}(a_0)\) and \(\varphi^n(a_0)\) is a \(K\)-quasigeodesic in the curve graph. In particular, if \((c_i)_i\) is such a geodesic and \(Z\) a subsurface such that \(\pi_Z(c_i)\neq \emptyset\) for every \(i\) between \(i_0\) and \(j_0\),
\[\dist_Z(\pi_Z(c_{i_0}),\pi_Z(c_{j_0}))\leq B.\]
\end{enumerate}
\end{lemma}

\begin{proof}
  \begin{enumerate}
  \item First note that there is a sequence $m_i\to \infty$, so that
    the geodesic representatives of $\varphi^{m_i}(a_0)$ Hausdorff
    converge to a lamination containing $\lambda^+$ (since $\lambda^+$
    is the attracting fixed point of $\varphi$ on the sphere of
    projective laminations). This implies that for any $Z$ there is a
    power $m$ so that
    \[ \dist_Z(\pi_Z(\varphi^m(a_0)), \pi_Z(\lambda^+)) \leq 4. \] 
      Thus it suffices to show that
    \[\dist_Z(\pi_Z(\varphi^n(a_0)),\pi_Z(\varphi^m(a_0)))\leq A\]
    for all $n,m$. This is a well-known fact, we sketch a proof for
    the convenience of the reader. First, extend the
    orbit $\varphi^k(a_0), k\in\mathbb{Z}$ to a $\varphi$--invariant
    path $\rho$. Since the orbit is a quasigeodesic, we may assume
    that $\rho$ is a $K$-quasigeodesic for some $K>0$.

	Let $B$ denote the constant from the Bounded Geodesic Image Theorem (Corollary~\ref{cor:qBGIT}) for $K$-quasigeodesics.
    Choosing $n_0 > 0$ large enough we see that if
    $\dist_Z(\pi_Z(\varphi^n(a_0)),\pi_Z(\varphi^m(a_0)))> A$, then
    there is an index $i$ so that
    $\dist_Z(\pi_Z(\varphi^i(a_0)),\pi_Z(\varphi^{i+n_0}(a_0)))>
    A-2B$. Namely, the Bounded Quasigeodesic Image Theorem
    %(Corollary~\ref{cor:qBGIT}) 
    allows us to discard initial and
    terminal segments outside the $2$--neighbourhood of $\partial Z$ where the path may 	have empty projection in $Z$, at the cost of $B$ each.

    It thus suffices to find a bound in the case where 
    $m=n+n_0$. But this follows immediately from the fact that the
    subsurface projection between $\varphi^n(a_0), \varphi^{n+n_0}(a_0)$ is always
    bounded by the intersection number of those curves, which equals the intersection 	
    number between $\varphi(a_0), \varphi^{n_0}(a_0)$, independent of $Z$.
  \item Suppose that $(c_i)_i$ is a geodesic in the surviving
    curve graph between \(\varphi^{-n}(a_0)\) and
    \(\varphi^n(a_0)\). Note that $(\varphi^j(a_0))_{j=-n}^n$ is
    a $k_0$-quasigeodesic between the same endpoints, where $k_0$
    depends only on $\varphi$ and $a_0$, since it acts loxodromically on the
    fine curve graph. By the Morse lemma, this implies that for each
    $i$ there is a power $j(i)$ so that
    \[ \dist(c_i, \varphi^{j(i)}(a_0)) \leq M. \]
    Using that $c_i$ is a geodesic, we obtain
    \[ |k - l| = \dist(c_k, c_l) \leq
      d(\varphi^{j(k)}(a_0), \varphi^{j(l)}(a_0)) + 2M \leq
      k_0|j(k)-j(l)| + k_0 + 2M. \] Now, $\varphi$ also acts on the
    curve graph as a loxodromic element, and so
    $(\varphi^j(a_0))_{j=-n}^n$ is a $k_1$-quasigeodesic in the
    curve graph. We thus have
    \[ \dist_\C(c_k,c_l) \geq \dist_\C(\varphi^{j(k)}(a_0),
      \varphi^{j(l)}(a_0)) - 2M \geq \frac{1}{k_1}|j(k)-j(l)| - k_1 -
      2M. \] Combining those two via $d(c_k,c_l) \geq d_\C(c_k,c_l)$, we see that $c_i$ is a
    \(\C(S)\)-quasigeodesic with constants depending only on $k_0,k_1,M$.

    The final sentence is now an immediate consequence of the Bounded
    Quasigeodesic Image Theorem (Corollary~\ref{cor:qBGIT}).
  \end{enumerate}
\end{proof}

\end{document}
